% preamble_zh.tex — shared packages, macros, and theorem environments for the
% CHINESE edition. Compile the master main_zh.tex with XeLaTeX.
%
% CJK is provided by xeCJK + the bundled Fandol fonts on a plain `book` class
% (NOT ctexbook). ctexbook clobbers cleveref's \label hook at \begin{document},
% which breaks every \cref (the notes use \cref heavily); book+xeCJK keeps
% cleveref fully working. Chinese structural localization (chapter headings,
% \contentsname, figure/table names) is supplied explicitly below.
% The pdfLaTeX-only font lines (inputenc/fontenc/lmodern) are intentionally omitted.

\usepackage[margin=1in]{geometry}
\usepackage{amsmath,amssymb,amsthm,mathtools}
\usepackage{braket}          % Dirac notation: \ket, \bra, \braket
\usepackage{bm}
\usepackage{graphicx}
\usepackage{booktabs}
\usepackage{enumitem}
\usepackage{xcolor}
\usepackage{tikz}
\usepackage{quantikz}
\usepackage[most]{tcolorbox}

% --- CJK support (xeCJK + Fandol) ---
\usepackage{xeCJK}
\setCJKmainfont{FandolSong-Regular.otf}[
  Path=/usr/local/texlive/2025/texmf-dist/fonts/opentype/public/fandol/,
  BoldFont=FandolHei-Regular.otf, ItalicFont=FandolKai-Regular.otf]
\setCJKsansfont{FandolHei-Regular.otf}[
  Path=/usr/local/texlive/2025/texmf-dist/fonts/opentype/public/fandol/]
\linespread{1.3}             % a little extra leading for CJK readability

\usepackage[colorlinks=true,linkcolor=blue,citecolor=blue,urlcolor=blue]{hyperref}
\usepackage[capitalize,noabbrev]{cleveref}

% --- Chinese structural localization (replaces what ctexbook would provide) ---
\renewcommand{\contentsname}{目录}
\renewcommand{\listfigurename}{插图目录}
\renewcommand{\listtablename}{表格目录}
\renewcommand{\bibname}{参考文献}
\renewcommand{\indexname}{索引}
\renewcommand{\figurename}{图}
\renewcommand{\tablename}{表}
\renewcommand{\appendixname}{附录}
\makeatletter
\renewcommand{\@makechapterhead}[1]{%
  \vspace*{40\p@}{\parindent\z@\raggedright\normalfont
    \ifnum\c@secnumdepth>\m@ne \huge\bfseries 第\,\thechapter\,章\par\nobreak\vskip 18\p@\fi
    \Huge\bfseries #1\par\nobreak\vskip 30\p@}}
\renewcommand{\@makeschapterhead}[1]{%
  \vspace*{40\p@}{\parindent\z@\raggedright\normalfont
    \Huge\bfseries #1\par\nobreak\vskip 30\p@}}
\makeatother

% ---------------------------------------------------------------------------
% Draft helpers
% ---------------------------------------------------------------------------
\newcommand{\todo}[1]{{\color{red}\textbf{[待办: #1]}}}
\newcommand{\note}[1]{{\color{purple}\small(#1)}}
\newcommand{\coverage}[1]{%
  \par\smallskip\noindent
  \textcolor{teal}{\textbf{本学期讲授（2026--2027 秋季学期）。}}\enspace
  {\small #1}\par\smallskip}

% ---------------------------------------------------------------------------
% Notation macros (identical to the English build)
% ---------------------------------------------------------------------------
\newcommand{\R}{\mathbb{R}}
\newcommand{\C}{\mathbb{C}}
\newcommand{\N}{\mathbb{N}}
\newcommand{\Z}{\mathbb{Z}}
\newcommand{\E}{\mathbb{E}}
\newcommand{\Prob}{\mathbb{P}}
\newcommand{\Hilb}{\mathcal{H}}
\newcommand{\ii}{\mathrm{i}}
\newcommand{\dd}{\mathrm{d}}
\newcommand{\ee}{\mathrm{e}}
\DeclareMathOperator{\tr}{tr}
\DeclareMathOperator{\Tr}{Tr}
\DeclareMathOperator{\Var}{Var}
\DeclareMathOperator{\Cov}{Cov}
\DeclareMathOperator{\rank}{rank}
\DeclareMathOperator{\diag}{diag}
\DeclareMathOperator{\poly}{poly}
\DeclareMathOperator{\polylog}{polylog}
\DeclareMathOperator*{\argmin}{arg\,min}
\DeclareMathOperator*{\argmax}{arg\,max}
\newcommand{\norm}[1]{\left\lVert #1 \right\rVert}
\newcommand{\abs}[1]{\left\lvert #1 \right\rvert}
\newcommand{\bigo}{\mathcal{O}}
\newcommand{\id}{\mathbb{I}}
\newcommand{\dg}{^{\dagger}}
\let\proj\relax
\newcommand{\proj}[1]{\ket{#1}\!\bra{#1}}

% Pauli / common gates
\newcommand{\Xgate}{\mathsf{X}}
\newcommand{\Ygate}{\mathsf{Y}}
\newcommand{\Zgate}{\mathsf{Z}}
\newcommand{\Hgate}{\mathsf{H}}
\newcommand{\CNOT}{\mathrm{CNOT}}
\newcommand{\Uf}{U_{f}}
\newcommand{\QFT}{\mathrm{QFT}}
\newcommand{\defeq}{\vcentcolon=}

% ---------------------------------------------------------------------------
% Theorem environments (Chinese names)
% ---------------------------------------------------------------------------
\theoremstyle{plain}
\newtheorem{theorem}{定理}[chapter]
\newtheorem{lemma}[theorem]{引理}
\newtheorem{proposition}[theorem]{命题}
\newtheorem{corollary}[theorem]{推论}
\theoremstyle{definition}
\newtheorem{definition}[theorem]{定义}
\newtheorem{example}[theorem]{例}
\newtheorem{exercise}{习题}[chapter]
\theoremstyle{remark}
\newtheorem{remark}[theorem]{注}
\renewcommand{\proofname}{证明}
\newenvironment{solution}{\begin{proof}[\normalfont\itshape 解.]}{\end{proof}}

% cleveref: Chinese cross-reference names
\crefname{theorem}{定理}{定理}      \Crefname{theorem}{定理}{定理}
\crefname{lemma}{引理}{引理}        \Crefname{lemma}{引理}{引理}
\crefname{proposition}{命题}{命题}  \Crefname{proposition}{命题}{命题}
\crefname{corollary}{推论}{推论}    \Crefname{corollary}{推论}{推论}
\crefname{definition}{定义}{定义}   \Crefname{definition}{定义}{定义}
\crefname{example}{例}{例}          \Crefname{example}{例}{例}
\crefname{remark}{注}{注}           \Crefname{remark}{注}{注}
\crefname{exercise}{习题}{习题}     \Crefname{exercise}{习题}{习题}
\crefname{chapter}{章}{章}          \Crefname{chapter}{章}{章}
\crefname{section}{节}{节}          \Crefname{section}{节}{节}
\crefname{equation}{式}{式}         \Crefname{equation}{式}{式}
\crefname{figure}{图}{图}           \Crefname{figure}{图}{图}
\crefname{table}{表}{表}            \Crefname{table}{表}{表}

% ---------------------------------------------------------------------------
% Per-lecture structural helpers (Chinese)
% ---------------------------------------------------------------------------
\newtcolorbox{objectives}{
  colback=blue!4, colframe=blue!45!black, title={学习目标},
  fonttitle=\bfseries, boxrule=0.5pt, arc=2pt, left=6pt, right=6pt, top=4pt, bottom=4pt}
\newcommand{\exercisesection}{\section*{习题}\addcontentsline{toc}{section}{习题}}
\newcommand{\readingsection}{\section*{延伸阅读}\addcontentsline{toc}{section}{延伸阅读}}
